\documentclass[letter,11pt,oneside]{article}
%%% HEREHEREHERE
%%% APPENDIX
%%% (insert (format "\n%s\n" (buffer-file-name)))
%%% (occur "\\(\\\\[a-z]*section\\|appendix\\|input\\|\\<include\\>\\)")

%%\documentclass[11pt,twocolumn]{article}
%%\usepackage[inline]{asymptote}       %% Inline asymptote diagrams
%%\usepackage{wglatex}                 %% Use this one and kill others.
\usepackage{color}                     %% colored letters {\color{red}{{text}}
\usepackage{fancyhdr}                  %% headers/footers
\usepackage{datetime}                  %% pick up tex date time 
\usepackage{lastpage}                  %% support page of ...lastpage
\usepackage{times}                     %% native times roman fonts
\usepackage{textcomp}                  %% trademark
\usepackage{amssymb,amsmath}           %% greek alphabet
\usepackage{parskip}                   %% blank lines between paragraphs, no indent
\usepackage{shortvrb}                  %% short verb use for tables
\usepackage{lscape}                    %% landscape for tables.
\usepackage{longtable}                 %% permit tables to span pages wg-longtable
\usepackage{multicol}                  %% Enhance footnotes/endnotes
\usepackage{url}                       %% Make URLs uniform and links in PDFs
\usepackage{enumerate}                 %% Allow letters/decorations for enumerations
\usepackage{endnotes}                  %% Enhance footnotes/endnotes
\usepackage{listings}                  %% Make URLs uniform and links in PDFs
\pdfadjustspacing=1                    %% force LaTeX-like character spacing
\usepackage{geometry}                  %% allow margins to be relaxed
%%\usepackage{wrapfig}                 %% permit wrapping figures.
%%\usepackage{subfigure}               %% images side by side.
\geometry{margin=1in}                  %% Allow narrower margins etc.
\usepackage[T1]{fontenc}               %% Better Verbatim Font.
\renewcommand*\ttdefault{txtt}        %% 
\usepackage[colorlinks=true]{hyperref} %% Make huperlinks within a PDF
\usepackage{natbib}                    %% bibitems
\usepackage{upquote}                   %% make programmer's quoetes in verbatim sections.
\usepackage{verbatim}
%% include background image (wg-document-page-background) 

\usepackage{graphicx}            %% Include pictures into a document
%% (wg-texdoc-inserttikz)


\def\documentisdraft{NOTDRAFT}

%% (wg-texdoc-isdraft)
%% (wg-texdoc-insert-fancy-headers)

%%\usepackage[bookmarks]{hyperref} %% Make huperlinks within a PDF
%%\usepackage{makeidx}             %% Make an index uncomment following line
%%\makeindex                       %%.. yeah this one, too. index{key} in text
%%



\definecolor{verbcolor}{rgb}{0.6,0,0}
\definecolor{darkgreen}{rgb}{0,0.4,0}
\newcommand\debate[1]{\textcolor{darkgreen}{DEBATE: #1} \marginpar{\textcolor{red}{DEBATE} }}
\newcommand{\ltodo}[2]{\marginpar{\textcolor{red}{ACTION: #1}\endnote{#2}}}
\renewcommand{\thefigure}{\thesection-\arabic{figure}}
\newcommand{\menu}{\ensuremath{\;\rightarrow\;}}
\newcommand{\dhl}[1]{{\color{verbcolor}{\texttt#1}}}
\definecolor{wglightgreen}{rgb}{0.88, 0.58, 0.88}
\newcommand{\wgtextbox}[1]{\noindent\fcolorbox{darkgreen}{wglightgreen}{%
    \minipage[t]{\dimexpr0.80\linewidth-2\fboxsep-2\fboxrule\relax}
        {#1}
    \endminipage}}

%%(wg-latex-snippet)

%%(wg-add-inline-images)  %% add inline images to the mix




%%Begin User Definitions: Hint: ~/.latex.defs and  latex.defs  
%%End User Definitions:

%% (wg-texdoc-adjust-paper-width)
%% (wg-texdoc-insert-hypersetup)
%% (wg-latex-tablet-page)


%%%%%%%%%%%%%%%%%%%%%%%%%%%%%%%%%%%%%%%%%%%%%%%%%%%%%%%%%%%%%%%%%%%%%%%%%%%%%


\begin{document}


%% (wg-latex-pretty-title-page)
%%%%%%%%%%%%%%%%%%%%%%%%%%%% START PRETTY TITLE PAGE %%%%%%%%%%%%%%%%%%%%%%%%%%%%%%%%%

\vskip 1.5cm


\begin{center}
 {\Huge PyRAF and Docker \\
~ \\
 {\Large Using Docker Containers for PyRAF Reduction} }
\end{center}

\vskip 2.5cm


%%%%%%%%%%%%%%%%%%%%%%%%%%%% END PRETTY TITLE PAGE %%%%%%%%%%%%%%%%%%%%%%%%%%%%%%%%%

%% (wg-texdoc-titleblock)

\setcounter{section}{0}
\pagenumbering{arabic}

\ifx\documentisdraft\drafttest
\linenumbers    %%%%%%%%%%%%% DRAFT
\fi

\section{Overview}

Using Docker to host PyRAF brings the advantage of running a OS properly
matched and configured to the PyRAF requirements. Advances with the Docker
ecosystem allow for a \dhl{docker-compose.yml} file to build the image
and make all the provide the necessary connections between the data source,
machine resources and the running container. Augmented with a bit of bash
scripting, starting and using PyRAF becomes routine.

A little policy about data organization helps to automate reduction routine.
Adding additional Python code, ala a Pyraflogin3.py like file, and a reduce.pyraf
file really expedite getting down to business. Additional scripts may be used
to help with big tasks, like fixing filenames, trimming spectra and images, annotating
the headers all the details that promote a good finished product.

This document provides a set of notes for the task at hand.

\subsection{System Example}

A few conditions are established:

\begingroup \fontsize{10pt}{10pt}
\selectfont
%%\begin{Verbatim} [commandchars=\\\{\}]
\begin{verbatim} 
- Share a few observatory-machine directories under Win11.
- Create mount points, configure the base user-machine OS to support SMB fileshares.
- Create a known directory policy/structure.
- Install the Docker Desktop environment, using the free Docker account.
- Login into Docker, and get ready to go:
  -- This establishes the proper permissions for the Docker engine to access the main user
     system.
- Prepare (git and by hand) all the files you will use/share with the team.
\end{verbatim}
\endgroup
%% \end{Verbatim}

Data is taken using a telescope system under Windows11 (ASCOM usual toolsuites
that require Windows) at a site in Chile. Observations are made in several other
locations around the world: Arizona, Colorado in the US and Newcastle UK. The
network is implemented with Tailscale, the Remote Win11 machine us using TightVNC
together with Tailscale. RealVNC viewer may be used simultaneously by the users,
facilitated by a Zoom\texttrademark meeting.

Observations at the remote end drop their files into /Users/ourobstech/Desktop/FS1Today.
This does not change night to night. In the morning, the files are moved to
/Users/ourobstech/Desktop/SpectroData, the FS1Today file is cleared and prepared
for the next time. In the past we simply renamed the FS1Today directory -- but when
it was shared on a permenant basis we lost that ability. 

During an observation run, we remotely access fresh files locally at /mnt/FS1Today.
We can move/access previous runs at /mnt/SpectroData. Seamless from the point of
the casual user.

In this one case, SMB file sharing is implemented on the Main machine, using the
small bash script:

\begingroup \fontsize{8pt}{8pt}
\selectfont
%%\begin{Verbatim} [commandchars=\\\{\}]
\begin{verbatim} 
#!/bin/bash

# get the password
read -r -s -p "Enter your password: " PASSWORD
ipaddr=$1
if test "$ipaddr" == "" ; then 
  ipaddr="100.95.214.76"
fi
sudo mount -t cifs //"$ipaddr"/SpectroData /mnt/SpectroData -o username=ourobstech,password="$PASSWORD",uid=1000,gid=1000
sudo mount -t cifs //"$ipaddr"/FS1Today /mnt/FS1Today       -o username=ourobstech,password="$PASSWORD",uid=1000,gid=1000
\end{verbatim}
\endgroup
%% \end{Verbatim}


An Intel I9-14000 Ubuntu 22.04 machine is running Docker. The PyRAF container
environmnet uses Ubuntu 24.04. D

\appendix
\renewcommand \thesection{\Alph{section}}

\section{SSH Keys}

files are in ~/.ssh

you get a public key, NAME sans extension, and the public key
NAME.pub

Never show anyone the private key.



zip that up, and save to the side.



\subsection{Backup}
\begingroup \fontsize{10pt}{10pt}
\selectfont
%%\begin{Verbatim} [commandchars=\\\{\}]
\begin{verbatim} 
https://www.youtube.com/watch?v=GxRu35fy-oY
~/.ssh

ed25519 cryptographically better -C "GAO", shorter and supported on all linux machines.

> gao_key
> passphrase

# Remember the passphrase (it is the 'password' to get machine via ssh. Not
# the actual password of the remote machine.

ssh-keygen -t ed25519 -f /home/$USER/.ssh/gao -q -N "sky is the limit" <<< y
# makes files gao and gao.pub
ssh-keygen -t ed25519 -f /home/$USER/.ssh/george -q -N "Jane! Stop this crazy thing!" <<< y




passphrase, increases strength of the key. locks the key.



eval $(ssh-agent)  # run and keep alive.


ssh-add <passphrase>  # allow direct remote login

Make a config file

Host gao
  Hostname <ipaddress>
  IdentityFile ~/.ssh/gao_key



\end{verbatim}
\endgroup
%% \end{Verbatim}

\section{Docker Configuration File}



\begingroup \fontsize{8pt}{8pt}
\selectfont
\verbatiminput{txt/docker-compose.yml}
\endgroup

%% use a bibitem approach to the references publications etc.
%% (wg-bibitem)



%%\begin{thebibliography}{80}
%%\usepackage{natbib}   %% bibitems
%%\end{thebibliography}

%%\clearpage
%%\addcontentsline{toc}{section}{Index}
%%\printindex %% www.cs.usask.ca/resources/tutorials/latex/notes/toc-index.pdf

% /home/git/external/BonMots/tex/Docker/Docker.tex

%% (wg-texdoc-endnotes)
\end{document}
